\section{Introduction}

E-graphs provide a compact representation of many equivalent expressions, enabling equality saturation to explore alternative rewrites without discarding earlier forms and to defer the selection of a preferred expression to an application-specific cost model~\cite{ogEgraphs1,egraphsOG,eqsat,eqsat-lmcs,egg}. %\Fix{Can we state some benefits / reasons why people would use e-graphs here? Basically provide some high level motivation for working on e-graphs at large before diving into limitations we address thereof.} Two
Orthogonal extensions complicate those keys: algebraic operators benefit
from sequence-, bag-, or set-valued children rather than explicit
associativity, commutativity, and idempotency (ACI) saturation; open terms
benefit from interfaces quotiented by consistent renaming.  Neither dimension
alone covers the case where an algebraic container holds a typed invocation
of an open e-class.

We give a proof-carrying composition of the typed slots and quotient space over the aforementioned algebraic properties in equality saturation. %\Fix{Of what? is what I am left thinking. Maybe some transition lead in is needed or just a clarification}.  
Each operator declares ports generated
by \texttt{One}, $K^J$ for $K\in\{\Seq,\Bag,\Set\}$, unary binding, and
binder blocks.  A child invocation $m*a$ embeds the exposed slots of e-class
$a$ into its caller.  The carrier determines only the certified sibling
quotient; a separate, operator- and path-specific certificate licenses
recursive same-head flattening.  Equality-certified permutations act on
exposed slots, while descriptor-certified automorphisms act locally inside a
binder block.  Rebuilding normalizes these actions jointly with union-find
leaders in the intended design without inferring child equality from parent
equality.  The checked trace theorem instead consumes local obligation
evidence and does not verify an operational rebuild implementation.

The proposed construction is quotient-first.  Leader finding may contract
support through a proper embedding, so canonicalization first extracts an
exact-support leader kernel and retains a source-to-kernel certificate. Such process 
%\Fix{Instead of `it' can we name the thing directly? It can be ambiguous so its good practice to avoid using it in academia writing to keep statements as clear as possible. }
then minimizes local leader and binder orbits before aggregating bags or
deduplicating sets, and only afterward minimizes the finite free-slot-renaming
orbit.  The checked mathematical layer proves exactness for any supplied
finite quotient presentation and validates supplied extraction and collision
records; showing that the concrete procedure constructs those records is a
separate refinement obligation.  A collision may create a parent edge only
when the composed certificate inhabits that edge's exact directed type.

We prove finite-unfolding soundness for abstract finite certified traces.
Their nine-step labels are classified as local obligation addition, rekeying,
composition, or removal, each carrying the needed typed endpoint evidence.
Under these premises, two certified unfoldings whose parent paths reach one
leader and whose caller maps align by a recorded symmetry word are equal in
every supplied model. %\Fix{I think it would be helpful to have some transition into this sentence, like`` To note,'' or something but not required} 
To note, the result does not establish concrete transition
semantics, unrestricted ACI completeness, or refinement of the Java code.

%\Fix{Transition framing for here, maybe something like} \Fix{To evaluate our e-graph representation on a concrete domain, }
To evaluate our e-graph construction on a concrete formal dataset, we explore an Alloy case study with seven related pipeline arms on a
frozen corpus and a targeted capability suite. The pipeline separates a targeted
binder-permutation capability from whole-pipeline structural consolidation;
the measurements are bounded implementation evidence, not a proof that
the adapter realizes the abstract construction. In our controlled benchmark, the slot-aware configurations recognized all 5,500 constructed equivalences, including all 500 binder-block permutations, compared with 22 for the De Bruijn baselines. Across 61,598 AST-distinct Alloy predicate pairs, the Certificate-Integrated IR recognized 4,088 correct student–oracle equivalences versus 2,160 for the De Bruijn baselines, with no observed zero-distance matches among pairs labeled incorrect.
%\Fix{I would add a sentence or two recapping the outcome of this with a key observation or two.}

Our contributions are: (i) a typed flexible-arity port grammar for slotted
invocations; (ii) a generic finite-orbit exactness theorem and certified
effective-support record interfaces with explicit parent-versus-interface
restriction; (iii) an abstract certificate-preserving trace contract and
finite-unfolding soundness theorem; and
(iv) a bounded Alloy study whose evidence boundary is stated separately from
the mathematical result.  Full statements and proof details appear in
Appendix~\ref{app:formal-details}. %the machine-checked correspondence is
%tracked by the \texttt{\string\leanref} tags attached to every formal unit.
